\begin{abstract}
A face verification system looks at two photos and decides if they show the same person. The problem is that a very small change to one photo, too small for a person to notice, can make the system give the wrong answer. Simple image filters such as JPEG compression, blurring and median filtering are a cheap way to remove such changes, but they tend to stop working once the attacker knows which filter is used. We tested this on the LFW dataset with a pretrained face model that we did not retrain. First, we found that the filters only help a little: at $\epsilon{=}4/255$ the best filter still lets PGD succeed 79\% of the time, and attackers who know the filter (BPDA) reach 90\% against JPEG-30 and 97\% against the median filter. Second, we built a detector that needs no training of the face model. For each pair of photos we apply a few cheap filters and check how much the similarity score and the face embedding change. Clean pairs change very little, attacked pairs change a lot. We feed these changes to simple classifiers from our course (LDA, Gaussian Bayes, SVM, boosted trees) and then use a Bayes minimum-risk rule to give one of three answers: Same, Different or Flag. The detector was trained only on FGSM and PGD, yet it brings the attack success rate of FGSM, PGD, CW and two adaptive attacks down to 2\% or less. The price is that \cleanflag\% of clean pairs also get flagged, and it does not catch our weak Square attack. These results should be read with care: our adaptive attacks are not the strongest possible, and everything comes from one seed, one dataset and one model.
\end{abstract}
